\documentclass[11pt]{article}
\usepackage[margin=1in]{geometry}
\usepackage{amsmath,amssymb,amsthm,mathtools}
\usepackage{graphicx}
\usepackage{xcolor}
\usepackage{microtype}
\usepackage{enumitem}
\usepackage[hidelinks]{hyperref}

\newtheorem{theorem}{Theorem}
\newtheorem{corollary}[theorem]{Corollary}
\newtheorem{proposition}[theorem]{Proposition}
\newtheorem{lemma}[theorem]{Lemma}
\theoremstyle{remark}
\newtheorem{remark}[theorem]{Remark}

\newcommand{\BH}{\mathcal B(\mathcal H)}
\newcommand{\JJ}{\mathcal J}
\newcommand{\trc}{\tau}
\newcommand{\deltaT}[2]{\Delta_{\trc}(#1,#2)}

\title{A regular-modulo-ideal subcase of the\\trace commutator problem}
\author{Candidate partial result for arXiv:1712.06702}
\date{11 August 2026}

\begin{document}
\maketitle

\begin{center}
\fbox{\parbox{0.91\textwidth}{\textbf{Status.}
Candidate partial result, likely valid.  The unrestricted question remains
open.  The proof below covers every pair for which either factor is von
Neumann regular modulo the trace ideal; in particular it covers every pair
with a closed-range factor.}}
\end{center}

\section{Source question}

Loreaux and Weiss ask the following as Question~1 of
\cite{LoreauxWeiss2018}.  Here $\JJ$ is a proper two-sided ideal of $\BH$ and
a trace is a linear functional $\trc:\JJ\to\mathbb C$ satisfying
$\trc(XS)=\trc(SX)$ for $X\in\JJ$ and $S\in\BH$.

\begin{quote}
Given a trace $\trc:\JJ\to\mathbb C$ and $A,B\in\BH$, does
\[
  \trc(AB)=\trc(BA) \qquad\text{whenever }AB,BA\in\JJ?
\]
\end{quote}

The statement appears on page~1 of the source paper:

\begin{center}
\includegraphics[width=\textwidth]{figures/open_problem_crop.png}
\end{center}

The source proves an affirmative answer for spectral traces and points to
nonspectral traces and quasinilpotent operators as possible sources of a
counterexample.  The theorem below gives a different, purely algebraic
affirmative subcase.

\section{Idea of the proof}

The key operation is to replace one factor $A$ by its defect from having an
inner inverse.  For an arbitrary bounded $C$, set $R=A-ACA$.  Although neither
$A$ nor $B$ need lie in the ideal, a pair of legal cyclic rotations shows that
the trace defect is preserved exactly:
\[
  \trc(AB)-\trc(BA)=\trc(RB)-\trc(BR).
\]
If the coset of $A$ is von Neumann regular modulo $\JJ$, then $C$ can be
chosen with $R\in\JJ$, and the right side vanishes by the defining trace
property.  For a closed-range operator one may take a bounded generalized
inverse and obtain $R=0$.

Two reductions identify what is left.  Matrix amplification turns any
counterexample into one in which both factors are square-zero.  Conversely,
after polar reduction to $A\geq0$, truncated generalized inverses replace $A$
by an arbitrarily small spectral tail without changing the trace defect.
Thus the unresolved mechanism is genuinely tied to nonclosed range and to the
lack of continuity of an arbitrary trace.

\section{The descent identity and the partial theorem}

We first state the argument in a unital algebra, since no Hilbert-space
geometry is used.

\begin{lemma}[Defect descent]
Let $\mathcal R$ be a unital complex algebra, let $\mathcal I\lhd\mathcal R$
be a two-sided ideal, and let $\trc:\mathcal I\to\mathbb C$ be a trace.  If
$a,b\in\mathcal R$ satisfy $ab,ba\in\mathcal I$, then for every
$c\in\mathcal R$, with $r=a-aca$, all of $rb,br$ lie in $\mathcal I$ and
\begin{equation}\label{eq:descent}
  \trc(ab)-\trc(ba)=\trc(rb)-\trc(br).
\end{equation}
\end{lemma}

\begin{proof}
The ideal property gives
\[
 rb=ab-acab\in\mathcal I,
 \qquad
 br=ba-baca\in\mathcal I.
\]
Moreover $cab\in\mathcal I$ and $ba\in\mathcal I$.  Two applications of the
trace identity, each with one factor in $\mathcal I$, give
\[
 \trc(acab)=\trc(caba)=\trc(baca).
\]
Expanding $rb-br$ and cancelling these equal middle terms proves
\eqref{eq:descent}.
\end{proof}

Recall that $a+\mathcal I$ is \emph{von Neumann regular} in
$\mathcal R/\mathcal I$ if there is $c\in\mathcal R$ such that
\[
 a-aca\in\mathcal I.
\]

\begin{theorem}[Regular-modulo-ideal trace commutation]\label{thm:regular}
Under the hypotheses of the lemma, if either $a+\mathcal I$ or
$b+\mathcal I$ is von Neumann regular in $\mathcal R/\mathcal I$, then
\[
 \trc(ab)=\trc(ba).
\]
\end{theorem}

\begin{proof}
Suppose $a-aca=r\in\mathcal I$.  Since $r\in\mathcal I$ and $b\in\mathcal R$,
$\trc(rb)=\trc(br)$.  Equation~\eqref{eq:descent} now gives the conclusion.
The case in which $b+\mathcal I$ is regular follows after interchanging
$a$ and $b$.
\end{proof}

\begin{corollary}[Closed-range factor]\label{cor:closed-range}
Let $\JJ$ be a proper $\BH$-ideal and let $\trc$ be any trace on $\JJ$.
Suppose $A,B\in\BH$ and $AB,BA\in\JJ$.  If either $A$ or $B$ has closed
range, then
\[
 \trc(AB)=\trc(BA).
\]
In particular this holds if either factor is Fredholm, semi-Fredholm, a
projection, a partial isometry, or bounded below on the orthogonal complement
of its kernel.
\end{corollary}

\begin{proof}
If $A$ has closed range, its Moore--Penrose inverse $C$ is bounded and
$ACA=A$.  Thus $A$ is von Neumann regular already in $\BH$, and
Theorem~\ref{thm:regular} applies.  The argument is symmetric in $A,B$.
\end{proof}

\section{Two exact reductions of the remaining problem}

\begin{proposition}[Square-zero reduction]\label{prop:squarezero}
The universal Loreaux--Weiss question is equivalent to its restriction to
pairs $P,Q$ satisfying $P^2=Q^2=0$.
\end{proposition}

\begin{proof}
Only one direction needs proof.  Suppose $A,B,\JJ,\trc$ form a counterexample.
On $\mathcal H\oplus\mathcal H$ define
\[
 P=\begin{pmatrix}0&A\\0&0\end{pmatrix},
 \qquad
 Q=\begin{pmatrix}0&0\\B&0\end{pmatrix}.
\]
Then $P^2=Q^2=0$ and
\[
 PQ=\begin{pmatrix}AB&0\\0&0\end{pmatrix},
 \qquad
 QP=\begin{pmatrix}0&0\\0&BA\end{pmatrix}
 \in M_2(\JJ).
\]
The functional
\[
 \trc_2((X_{ij}))=\trc(X_{11})+\trc(X_{22})
\]
is a trace on $M_2(\JJ)$.  Therefore
$\trc_2(PQ)=\trc(AB)\neq\trc(BA)=\trc_2(QP)$.  Since
$\mathcal H\oplus\mathcal H$ is again a separable infinite-dimensional
Hilbert space, this is a counterexample of the required square-zero form.
\end{proof}

Combining Proposition~\ref{prop:squarezero} with
Corollary~\ref{cor:closed-range} shows that a negative answer, if one exists,
already exists among square-zero factors whose ranges are both nonclosed.

\begin{proposition}[Positive spectral-tail localization]\label{prop:tail}
Suppose $A\geq0$, $B\in\BH$, and $AB,BA\in\JJ$.  For every $\epsilon>0$ let
\[
 E_\epsilon=\mathbf 1_{[0,\epsilon)}(A),
 \qquad R_\epsilon=AE_\epsilon.
\]
Then $R_\epsilon B,BR_\epsilon\in\JJ$, $\|R_\epsilon\|\leq\epsilon$, and
\[
 \deltaT{A}{B}:=\trc(AB)-\trc(BA)
 =\trc(R_\epsilon B)-\trc(BR_\epsilon).
\]
Moreover, every putative counterexample can first be reduced to one with a
positive first factor.
\end{proposition}

\begin{proof}
Define the bounded Borel function
\[
 f_\epsilon(t)=
 \begin{cases}t^{-1},&t\geq\epsilon,\\0,&0\leq t<\epsilon,
 \end{cases}
 \qquad C_\epsilon=f_\epsilon(A).
\]
Functional calculus gives
$A-AC_\epsilon A=A\mathbf 1_{[0,\epsilon)}(A)=R_\epsilon$.
The descent lemma gives the trace identity and ideal membership, and the norm
bound is immediate.

For the final assertion write the polar decomposition $A=U|A|$.  Since
$|A|B=U^*AB\in\JJ$, cyclicity gives
\[
 \trc(AB)=\trc(U|A|B)=\trc(|A|BU).
\]
Also $|A|BU=U^*ABU\in\JJ$ and $BU|A|=BA\in\JJ$.  Hence the trace defect of
$(A,B)$ equals that of the pair $(|A|,BU)$, whose first factor is positive.
\end{proof}

\section{Verification, limitations, and novelty status}

The proof uses only finite algebraic expansions and legal applications of
$\trc(XS)=\trc(SX)$ with $X\in\JJ$.  In particular, it assumes no positivity,
normality, or continuity of the trace.  The two identities a verifier should
check first are
\[
 \trc(acab)=\trc(caba)=\trc(baca)
\]
and the diagonal-block values of $PQ$ and $QP$ in
Proposition~\ref{prop:squarezero}.  Both have been checked directly.

This is a partial result only.  Theorem~\ref{thm:regular} says nothing when
both cosets are nonregular in $\BH/\JJ$.  Proposition~\ref{prop:tail} does not
permit a limit as $\epsilon\downarrow0$, because an arbitrary trace on an
algebraic operator ideal may be completely discontinuous.  The nonspectral
quasinilpotent construction of Dykema--Kalton \cite{DykemaKalton1998}, noted by
the source authors as a possible counterexample route, still lacks a
factorization whose two products remain in the same ideal.

A bounded search on 11 August 2026 used the exact title and question, the
authors' later papers, the phrases ``$AB,BA$ in an ideal'', ``trace commutator
property'', ``closed range'', and ``von Neumann regular modulo an ideal'', and
the related arXiv papers 2103.09215 and 2203.11995.  It found no resolution of
Question~1 and no occurrence of the regular-modulo-ideal theorem above.  This
is not a claim of definitive novelty; novelty confidence is moderate and a
specialist literature review is recommended.

\paragraph{Human-review recommendation.}
Verify the two cyclic rotations in the descent lemma and the amplification
trace in Proposition~\ref{prop:squarezero}.  If accepted, retain this as a
substantial partial theorem and structural reduction, not as a full solution.

\begin{thebibliography}{9}
\bibitem{LoreauxWeiss2018}
J. Loreaux and G. Weiss,
\emph{Traces on ideals and the commutator property},
arXiv:1712.06702, 2017.

\bibitem{DykemaKalton1998}
K. Dykema and N. J. Kalton,
\emph{Spectral characterization of sums of commutators. II},
J. Reine Angew. Math. \textbf{504} (1998), 127--137;
arXiv:math/9709208.
\end{thebibliography}

\end{document}

